\documentclass[aspectratio=169,10pt]{beamer}
\usetheme{Madrid}
\usecolortheme{default}

\usepackage{hyperref}
\usepackage{tikz}
\usepackage{amsmath,amssymb}

% --- Custom Colors ---
\definecolor{unibblue}{RGB}{0,51,102}
\definecolor{unibgold}{RGB}{212,175,55}
\definecolor{unibgreen}{RGB}{0,128,0}

\setbeamercolor{palette primary}{bg=unibblue,fg=white}
\setbeamercolor{palette secondary}{bg=unibblue!85,fg=white}
\setbeamercolor{palette tertiary}{bg=unibblue!70,fg=white}
\setbeamercolor{structure}{fg=unibblue}
\setbeamercolor{block title}{bg=unibblue!85,fg=white}
\setbeamercolor{block body}{bg=unibblue!10,fg=black}
\setbeamercolor{alertblock title}{bg=unibgold!90,fg=black}
\setbeamercolor{alertblock body}{bg=unibgold!15,fg=black}

\title{Persamaan Diferensial}
\subtitle{Minggu 11: Persamaan Gelombang 1D pada Saluran Transmisi Elektromagnetik}
\author{Novalio Daratha \& Muhammad Arfan}
\institute{Program Studi Teknik Elektro\\Universitas Bengkulu}
\date[\today]{2026 \\ \vspace{2mm}\small\textit{Tanggal Kompilasi: \today}}

\begin{document}

\begin{frame}
  \titlepage
\end{frame}

\begin{frame}{Capaian Pembelajaran (CPMK 4)}
  \begin{itemize}
    \item \textbf{CPMK 4.1:} Menurunkan Persamaan Gelombang 1D dari persamaan Telegraf pada saluran transmisi.
    \item \textbf{CPMK 4.2:} Menyelesaikan persamaan gelombang menggunakan Solusi Perjalanan Gelombang d'Alembert.
    \item \textbf{CPMK 4.3:} Menganalisis gelombang berdiri dan fenomena refleksi pada ujung saluran.
    \item \textbf{CPMK 4.4:} Mengamati simulasi numerik perambatan gelombang FDTD menggunakan Julia.
  \end{itemize}
\end{frame}

\begin{frame}{Daftar Isi}
  \tableofcontents
\end{frame}

\section{Penurunan dari Rangkaian Saluran Transmisi}
\begin{frame}{Persamaan Telegraf (Lossless Line)}
  Tinjau segmen saluran transmisi mikro sepanjang $\Delta x$ dengan induktansi $L'$ dan kapasitansi $C'$ per meter:
  $$ \frac{\partial v}{\partial x} = -L' \frac{\partial i}{\partial t} \quad \text{(Hukum Faraday/KVL)} $$
  $$ \frac{\partial i}{\partial x} = -C' \frac{\partial v}{\partial t} \quad \text{(Hukum Arus Maxwell/KCL)} $$
  \pause
  Turunkan persamaan tegangan terhadap $x$ dan persamaan arus terhadap $t$:
  $$ \frac{\partial^2 v}{\partial x^2} = -L' \frac{\partial^2 i}{\partial x \partial t} = -L' \frac{\partial}{\partial t}\left(-C'\frac{\partial v}{\partial t}\right) = L'C' \frac{\partial^2 v}{\partial t^2} $$
  \pause
  \begin{alertblock}{Persamaan Gelombang 1D Standar}
    $$ {\color{unibblue} \frac{\partial^2 v}{\partial t^2} = c^2 \frac{\partial^2 v}{\partial x^2}}, \quad c = \frac{1}{\sqrt{L'C'}} \quad \text{(Kecepatan Rambat Gelombang)} $$
  \end{alertblock}
\end{frame}

\section{Solusi d'Alembert: Gelombang Berjalan}
\begin{frame}{Formulasi Solusi d'Alembert}
  \begin{block}{Teorema d'Alembert}
    Solusi umum dari $\frac{\partial^2 u}{\partial t^2} = c^2 \frac{\partial^2 u}{\partial x^2}$ adalah superposisi dari dua gelombang yang merambat berlawanan arah:
    $$ {\color{unibblue} u(x,t) = f(x - ct) + g(x + ct)} $$
    \begin{itemize}
      \item $f(x - ct)$: Gelombang datang (\textit{forward traveling wave}) merambat ke arah $+x$.
      \item $g(x + ct)$: Gelombang pantul (\textit{backward traveling wave}) merambat ke arah $-x$.
    \end{itemize}
  \end{block}
  \pause
  Jika diberikan profil awal $u(x,0) = \phi(x)$ dan kecepatan awal $\frac{\partial u}{\partial t}(x,0) = 0$:
  $$ u(x,t) = \frac{1}{2}\left[ \phi(x - ct) + \phi(x + ct) \right] $$
\end{frame}

\section{Refleksi Gelombang pada Ujung Saluran}
\begin{frame}{Kondisi Batas Saluran Transmisi}
  \begin{columns}
    \begin{column}{0.5\textwidth}
      \begin{block}{Ujung Hubung Singkat ($x = L$ Ground)}
        \begin{itemize}
          \item Tegangan $v(L,t) = 0$ (Batas Dirichlet).
          \item Koefisien Refleksi: $\Gamma = -1$.
          \item Gelombang pantul berbalik polaritas (pergeseran fase $180^\circ$).
        \end{itemize}
      \end{block}
    \end{column}
    \begin{column}{0.5\textwidth}
      \pause
      \begin{block}{Ujung Terbuka ($x = L$ Open)}
        \begin{itemize}
          \item Arus $i(L,t) = 0 \implies \frac{\partial v}{\partial x}(L,t) = 0$ (Batas Neumann).
          \item Koefisien Refleksi: $\Gamma = +1$.
          \item Gelombang pantul sefase (tegangan berlipat ganda).
        \end{itemize}
      \end{block}
    \end{column}
  \end{columns}
\end{frame}

\section{Praktikum Julia \& Kuis}
\begin{frame}{Simulasi Julia: Perambatan Gelombang 1D}
  \begin{block}{Metode FDTD (Finite Difference Time Domain)}
    $$ u_i^{n+1} = 2(1 - r^2)u_i^n + r^2(u_{i+1}^n + u_{i-1}^n) - u_i^{n-1}, \quad r = \frac{c \Delta t}{\Delta x} \le 1 $$
    \begin{itemize}
      \item Memvisualisasikan pemecahan pulsa Gauss menjadi 2 gelombang simetris d'Alembert.
      \item Menghasilkan diagram spatiotemporal *heatmap* 2D.
    \end{itemize}
  \end{block}
\end{frame}

\begin{frame}{Kuis Interaktif 11}
  \textbf{Soal:} Saluran kabel koaksial memiliki $L' = 0.25\ \mu\text{H/m}$ dan $C' = 100\text{ pF/m}$. Berapakah kecepatan rambat sinyal dan waktu tempuh untuk kabel sepanjang $300\text{ meter}$?
  \pause
  \begin{block}{Jawaban:}
    \begin{enumerate}
      \item $c = \frac{1}{\sqrt{0.25 \times 10^{-6} \times 100 \times 10^{-12}}} = \frac{1}{\sqrt{25 \times 10^{-18}}} = \frac{1}{5 \times 10^{-9}} = 2 \times 10^8\text{ m/s} = \frac{2}{3}c$ (kecepatan cahaya). \pause
      \item Waktu tempuh: $t = \frac{300\text{ m}}{2 \times 10^8\text{ m/s}} = 1.5\ \mu\text{s} = 1500\text{ ns}$.
    \end{enumerate}
  \end{block}
\end{frame}

\begin{frame}{Rangkuman}
  \begin{enumerate}
    \item Gelombang tegangan dan arus pada saluran transmisi merupakan konsekuensi langsung dari hukum Maxwell.
    \item Kecepatan rambat hanya bergantung pada karakteristik geometris dan material dielektrik ($L', C'$).
    \item Analisis d'Alembert dan FDTD menjadi instrumen esensial dalam desain transmisi data berkecepatan tinggi.
  \end{enumerate}
\end{frame}

\end{document}
